\pdfoutput=1
\documentclass[11pt]{amsart}
\usepackage{amsmath,amssymb,amsthm}
\usepackage[margin=1.2in]{geometry}
\usepackage{booktabs,adjustbox}
\usepackage{xcolor}
\usepackage[colorlinks=true,linkcolor=blue,citecolor=blue,urlcolor=blue]{hyperref}
\hypersetup{pdfauthor={Bernd Johannes Wuebben},
 pdftitle={Canonical metrics on singular moduli of K3 bundles},
 pdfsubject={Canonical Ricci-flat metrics and quantitative tangent cones on K3 bundle moduli},
 pdfkeywords={K3 surface; moduli of bundles; singular Ricci-flat metric; hyperkahler quotient; metric tangent cone}}
\setlength{\emergencystretch}{2em}
\newtheorem{theorem}{Theorem}[section]
\newtheorem{proposition}[theorem]{Proposition}
\newtheorem{lemma}[theorem]{Lemma}
\newtheorem{corollary}[theorem]{Corollary}
\theoremstyle{definition}
\newtheorem{definition}[theorem]{Definition}
\newtheorem{example}[theorem]{Example}
\theoremstyle{remark}
\newtheorem{remark}[theorem]{Remark}
\newcommand{\CC}{\mathbb C}
\newcommand{\RR}{\mathbb R}
\newcommand{\ZZ}{\mathbb Z}
\newcommand{\OO}{\mathcal O}
\newcommand{\End}{\operatorname{End}}
\newcommand{\Hom}{\operatorname{Hom}}
\newcommand{\Ext}{\operatorname{Ext}}
\newcommand{\Aut}{\operatorname{Aut}}
\newcommand{\tr}{\operatorname{tr}}
\newcommand{\rk}{\operatorname{rk}}
\newcommand{\Sing}{\operatorname{Sing}}
\newcommand{\Sym}{\operatorname{Sym}}
\newcommand{\ch}{\operatorname{ch}}
\newcommand{\GL}{\operatorname{GL}}
\newcommand{\SU}{\operatorname{SU}}
\newcommand{\U}{\operatorname{U}}
\newcommand{\reg}{\mathrm{reg}}
\newcommand{\str}{\mathrm{str}}
\newcommand{\WP}{\mathrm{WP}}
\newcommand{\can}{\mathrm{can}}
\newcommand{\ddc}{dd^c}
\newcommand{\dbar}{\bar\partial}
\newcommand{\Cone}{\mathcal C}
\newcommand{\hkr}{\mathbin{/\!\!/\!\!/}}
\title[Canonical metrics on singular K3 moduli]{Canonical metrics on singular moduli of K3 bundles}
\author{Bernd Johannes Wuebben}
\subjclass[2020]{14J60 (primary); 14J28, 32Q20, 53C26 (secondary)}
\keywords{K3 surface, moduli of bundles, singular Ricci-flat metric,
hyperk\"ahler quotient, metric tangent cone}
\date{November 8, 2026\\This version: November 8, 2026}

\begin{document}
\begin{abstract}
The canonical singular Ricci-flat metric on a class of compact K3
bundle moduli spaces admits quadratic approximation by its
hyperk\"ahler tangent cones in holomorphic coordinates. The spaces
parametrize twice a primitive Mukai vector on a K3 surface of Picard
rank one; a rank bound ensures that every semistable sheaf and every
slope Jordan--H\"older factor is locally free. These spaces have
positive-dimensional singular strata and admit neither a symplectic
resolution nor a smoothing. We identify the metric in an ample
determinant class with the harmonic $L^2$ metric, identify its tangent
cones as the flat hyperk\"ahler quotients of the extension spaces,
and prove that its regular-locus completion is the original moduli
space. The comparison is uniform through changes of stabilizer in
each fixed local chart. The proof combines quantitative local models
for Hermitian--Einstein metrics with exact quadraticity in harmonic
Kuranishi coordinates, continuous Weil--Petersson potentials, and a
determinant-class calculation.
\end{abstract}
\maketitle
\setcounter{tocdepth}{1}
%\tableofcontents
%\newpage
\section{Introduction}\label{sec:introduction}

A projective symplectic variety carries a unique singular Ricci-flat
K\"ahler metric in each K\"ahler class, in the sense of
Eyssidieux--Guedj--Zeriahi \cite[Theorem~7.5 and Corollary~7.6]{egz09}.
The existence theorem does not by itself describe the local metric
geometry at a singular point. In a moduli space of bundles, a second
description is available: Hermitian--Einstein connections define a
hyperk\"ahler $L^2$ metric on the stable locus. We identify these two
metrics on a family of singular compact K3 moduli spaces and obtain
a quadratic comparison with explicit hyperk\"ahler cones at every
singular point.

The simplest example is already a compact space with two
positive-dimensional singular strata. Let $S$ be a quartic K3 surface
with $\operatorname{Pic}(S)=\ZZ H$, $H^2=4$, and put
\begin{equation}\label{eq:quartic}
 v_0=(5,H,0),\qquad v=2v_0=(10,2H,0),\qquad M=M_H(v).
\end{equation}
Here $M_H(v)$ denotes the moduli space of $H$-Gieseker-semistable
sheaves of Mukai vector $v$. This space has complex dimension $18$.
Every semistable sheaf is a vector bundle, and its slope
Jordan--H\"older factors are vector bundles as well. If
$N=M_H(v_0)$, then $N$ is a smooth compact hyperk\"ahler sixfold,
and the singular locus and its diagonal are
\begin{equation}\label{eq:quartic-strata}
 \Sigma=\Sym^2N,\qquad \Delta=\{[G\oplus G]:[G]\in N\}.
\end{equation}
Thus $\dim_{\CC}\Sigma=12$ and $\dim_{\CC}\Delta=6$.
The variety $M$ has terminal, locally factorial symplectic
singularities, admits no symplectic resolution, and has no smoothing.
These assertions follow from the moduli results recalled in
Section~\ref{sec:moduli}. The exclusion of a non-locally-free
boundary is proved there by a Mukai-pairing inequality.

\subsection{The family and the main result}
Let $S$ be a projective K3 surface with
$\operatorname{Pic}(S)=\ZZ H$, where $H$ is ample. Write the Mukai
vector of a sheaf as $v(F)=(\rk F,c_1(F),\chi(F)-\rk F)$, with
pairing
\[
 \langle(r,c,s),(r',c',s')\rangle=c\cdot c'-rs'-r's.
\]
We impose the conditions
\begin{equation}\label{eq:family}
 v_0=(r_0,aH,s_0),\qquad
 \gcd(r_0,aH^2)=1,\qquad q:=v_0^2\ge4,
 \qquad r_0\ge q+1,
\end{equation}
and set
\begin{equation}\label{eq:MN}
 M=M_H(2v_0),\qquad N=M_H(v_0).
\end{equation}
The integer $q$ is even. Fix the Ricci-flat K\"ahler form
$\omega\in c_1(H)$ on $S$.
Proposition~\ref{prop:no-bubbling} proves that slope and Gieseker
semistability agree for the sheaves under consideration.
The moduli results in Section~\ref{sec:moduli} identify
the stable locus $M^s$ with the regular locus of $M$. The
Hermitian--Einstein metrics therefore define the harmonic $L^2$
metric $g^{L^2}$ on $M_{\reg}$.

Here is the local cone that occurs in the theorem. For a polystable
bundle $E$, let
\begin{equation}\label{eq:cone-intro}
 V_E=H^1(S,\End E),\qquad K_E=\Aut(E,h_E),\qquad
 \Cone_E=V_E\hkr_0K_E.
\end{equation}
Harmonic representatives give $V_E$ its flat hyperk\"ahler
structure. The last expression means the simultaneous zero set of
the three moment maps, divided by the unitary automorphism group;
an ineffective scalar circle can be removed. The regular part of
this quotient carries the reduction metric $g_{\Cone_E}$. Its full
length metric is denoted by $d_{\Cone_E}$, and
$r(x)=d_{\Cone_E}(0,x)$. Equivalently, $\Cone_E$ is the flat
hyperk\"ahler quotient of the doubled Ext-quiver of the stable
summands of $E$. Section~\ref{sec:harmonic} specifies the harmonic
and moment-map conventions.

\begin{theorem}\label{thm:main}
Assume \eqref{eq:family}. There is an ample Donaldson determinant
line bundle $\lambda$ on $M$ and a constant $\kappa>0$ with the
following properties.
\begin{enumerate}
\item The K\"ahler form of $g^{L^2}$ extends to the canonical
singular Ricci-flat current in the class $\kappa c_1(\lambda)$.
It has continuous local potentials on $M$.
\item At every polystable point $[E]\in M$ there is a
biholomorphism of germs
\[
 \Psi_E:(\Cone_E,0)\longrightarrow(M,[E])
\]
for which, on the regular locus of a sufficiently small
neighbourhood,
\begin{equation}\label{eq:main-comparison}
 (1-C_Er^2)g_{\Cone_E}
 \le \Psi_E^*g^{L^2}
 \le (1+C_Er^2)g_{\Cone_E}.
\end{equation}
The same inequality holds on tangent vectors to every smooth
orbit-type stratum, with its harmonic metric. Within this fixed
chart, $C_E$ is independent of stabilizer type and of the ratios
of positive orbit eigenvalues.
\end{enumerate}
\end{theorem}

The determinant convention, including its positive sign, is given
in Section~\ref{sec:canonical}. The constants and chart radii in
\eqref{eq:main-comparison} may depend on the fixed bundle $E$,
its harmonic dimensions and its positive background elliptic
gaps. The theorem does not assert one radius or comparison
constant for all possible centred charts.

\begin{corollary}\label{cor:completion-main}
The metric completion of $(M_{\reg},g^{L^2})$ is a compact geodesic
space canonically homeomorphic to $M$. At every $p=[E]\in M$ its
metric tangent cone is unique and is $(\Cone_E,d_{\Cone_E},0)$.
In the chart of Theorem~\ref{thm:main},
\begin{equation}\label{eq:radius-main}
 d(p,\Psi_E(x))=r(x)(1+O_E(r(x)^2)).
\end{equation}
The Gromov--Hausdorff distance between the closed unit balls in
$(M,\epsilon^{-1}d,p)$ and in $\Cone_E$ is $O_E(\epsilon^2)$.
\end{corollary}

The completion assertion requires a path argument in addition to
the tensor estimate. We prove that paths along singular strata can
be approximated in length by regular paths. The perturbation is
transverse to the stratum in a product slice; radial dilation alone
preserves the singular stratum.

In the example \eqref{eq:quartic}, the two singular tangent cones
are
\begin{equation}\label{eq:quartic-cones}
 \begin{aligned}
 \Cone_{G_1\oplus G_2}
  &\cong \CC^{12}\times(T^*\CC^4\hkr_0\U(1)),
       &&G_1\not\cong G_2,\\
 \Cone_{G\oplus G}
  &\cong \CC^6\times
       ((\CC^6\otimes\mathfrak{sl}_2)\hkr_0\SU(2)).
 \end{aligned}
\end{equation}
The first normal factor is the closure of the minimal nilpotent
orbit in $\mathfrak{sl}_4$, of complex dimension six. The second
normal factor has complex dimension twelve. These products are
metric products, with the flat factors normalized by the harmonic
inner products.

The approximation in \eqref{eq:main-comparison} is not an exact
local isometry. Proposition~\ref{prop:not-isometric} proves that
no singular point has a neighbourhood isometric to a neighbourhood
of the vertex in its tangent cone. The reason is the intrinsic
geometry of the singular strata: their induced length metrics
contain the nonflat hyperk\"ahler metric of $N$. This does not
assert that the exponent two is optimal.

\subsection{Relation to earlier work}
The algebraic local models for these singularities belong to the
deformation theory of sheaves on K3 surfaces. Kaledin--Lehn--Sorger
established the structure and nonexistence of symplectic resolutions
for the relevant nonprimitive moduli spaces \cite{kls06}.
Arbarello--Sacc\`a identified the local quiver germs under
a formality hypothesis \cite[Propositions~4.4 and~6.1]{as18},
and Budur--Zhang proved that hypothesis for polystable sheaves
on projective K3 surfaces \cite[Theorem~1.2]{bz19}.
Perego--Rapagnetta established the irreducible
symplectic-variety properties used here \cite{pr23}.

On the analytic side, Itoh proved that the $L^2$ metric on the
irreducible Einstein--Hermitian moduli space over a K3 surface is
hyperk\"ahler \cite[Theorem~2.1 and Remark~(iii)]{itoh88}.
Verbitsky's theory of hyperholomorphic and projectively
hyperholomorphic bundles gives the underlying quadraticity
mechanism \cite[Theorems~6.2, 7.1 and 11.2, arXiv version]{verbitsky96}.
We use a surface calculation to show that the obstruction map is
quadratic in the particular harmonic coordinates needed for the
metric estimate. An unspecified analytic change of coordinates
would not preserve the stated quadratic error.

The quantitative local comparison for the harmonic metric is
imported from \cite[Theorem~1.1(iii)]{wuebben15}. Its application
here uses the full harmonic Kuranishi locus, including repeated
stable factors. The finite-dimensional metric slices and
localization arguments are those of \cite[Sections~4--6]{wuebben14}.
We state the required versions below. The additional global step
uses continuous Weil--Petersson potentials from
Buchdahl--Schumacher \cite[Theorem~11]{buchdahl21}, the curvature
formula of Biswas--Schumacher \cite[Proposition~4.3]{biswas16},
and uniqueness of the singular Ricci-flat current. Continuous
Weil--Petersson potentials themselves are an established input.

Namikawa--Odaka study canonical metrics and conical structures on
symplectic singularities, including a theorem under a symplectic
resolution or smoothing hypothesis and higher-order asymptotics
at isolated singularities \cite[Theorem~1.1 and Corollary~6.8]{no26}.
Odaka's subsequent result gives a hyperk\"ahler cone structure on
every symplectic singularity germ \cite[Corollary~1.3]{odaka26}.
The distinction relevant here is between the existence of a local
cone metric and the asymptotics of the metric selected by an ample
class on a compact moduli space. Theorem~\ref{thm:main} concerns
the latter metric and gives a relative estimate across its
positive-dimensional singular strata. It makes no assertion about
higher derivatives near lower strata or about the complete
Donaldson--Sun degeneration data.

\subsection{Proof and organization}
The local argument begins with the canonical harmonic
reconstruction $a=x-\bar\partial^*G(a\wedge a)$.
On a Ricci-flat K3 surface, harmonic endomorphism-valued
two-forms have parallel coefficients. Integration by parts
then annihilates the harmonic projection of every product
containing the coexact correction. Thus the obstruction
is exactly quadratic in the coordinates in which the
Hermitian--Einstein comparison has quadratic error.
The holomorphic family and the metric reconstruction
classify the same complex gauge orbit, so the resulting
chart is simultaneously holomorphic and quantitative.
Transverse perturbations in product slices then identify
the full stratified length distance with the completion
of the regular-locus metric.

For the global identification, the rank inequality forces
the defect cycles of all slope graded objects to vanish.
The Donaldson--Uhlenbeck morphism is consequently finite,
making its determinant polarization ample. A discriminant
calculation identifies this class with the harmonic
K\"ahler class. Continuous local potentials, together
with a depth-based cohomology restriction theorem, produce
a global continuous potential in that class. The
hyperk\"ahler volume identity then places the resulting
current in the uniqueness theorem of
Eyssidieux--Guedj--Zeriahi.

Section~\ref{sec:moduli} proves that the rank bound excludes both
non-locally-free sheaves and non-locally-free graded factors, and
describes the two singular strata. Section~\ref{sec:harmonic}
proves exact harmonic quadraticity and applies the quantitative
Hermitian--Einstein comparison in the same holomorphic quotient
coordinates. Section~\ref{sec:distances} identifies the completion
distance and derives the rescaling estimates.
Section~\ref{sec:canonical} identifies the ample determinant class
and proves the global Ricci-flat metric statement.
Section~\ref{sec:strata} proves the failure of local isometry.

All dimensions below are complex unless otherwise indicated.
We use $d^c=\frac{i}{2}(\bar\partial-\partial)$, so
$\ddc=i\partial\bar\partial$. A smooth stratum carries its
harmonic $L^2$ metric; a distance on the entire stratified space
will always be specified by its allowed paths.

\section{A compact moduli space of bundles}\label{sec:moduli}

We first show that the rank condition in \eqref{eq:family} excludes
both singular sheaves and singular stable factors. The second
assertion is needed for the determinant polarization: the cycle
recorded by the Donaldson--Uhlenbeck morphism is the defect of the
slope graded object, not just the defect of the original sheaf.

\begin{proposition}\label{prop:no-bubbling}
Assume \eqref{eq:family}. Every slope-semistable torsion-free sheaf
of Mukai vector $2v_0$ is locally free. If it is strictly
slope-semistable, its slope Jordan--H\"older factors are two stable
bundles of vector $v_0$. Every slope-semistable sheaf of vector
$v_0$ is stable and locally free. For these two Mukai vectors,
slope and Gieseker semistability agree, as do stability and
S-equivalence.
\end{proposition}

\begin{proof}
We use two elementary properties of stable sheaves on a K3
surface. If $A$ is stable, simplicity and Serre duality give
\begin{equation}\label{eq:stable-bound}
 v(A)^2=\dim\Ext^1(A,A)-2\ge-2.
\end{equation}
If $A$ is slope-stable and torsion-free, $A^{**}$ is a stable
bundle, and
\[
 v(A^{**})=v(A)+(0,0,\operatorname{length}(A^{**}/A)).
\]
The reflexive hull of a slope-semistable sheaf is semistable.

Let $F$ have vector $2v_0$, and write
$\ell=\operatorname{length}(F^{**}/F)$. If $F^{**}$ is stable,
\eqref{eq:stable-bound} implies
\[
 4q-4r_0\ell\ge-2.
\]
Since $r_0\ge q+1$, this forces $\ell=0$.

Suppose instead that $F^{**}$ is strictly slope-semistable.
An equal-slope factor of rank $r'$ and first Chern class $a'H$
satisfies
\[
 r_0a'H^2=r'aH^2.
\]
Coprimality in \eqref{eq:family} forces $r_0\mid r'$.
There are therefore exactly two stable factors $A_1,A_2$, each
of rank $r_0$ and first Chern class $aH$. Put
\[
 v(A_i)=(r_0,aH,s_i),\qquad
 \ell_i=\operatorname{length}(A_i^{**}/A_i).
\]
Additivity gives $s_1+s_2=2s_0+\ell$. Applied to $A_i^{**}$,
\eqref{eq:stable-bound} gives
\begin{equation}\label{eq:factor-inequality}
 q-2r_0(s_i+\ell_i-s_0)\ge-2.
\end{equation}
Because $2r_0>q+2$, integrality implies
$s_i+\ell_i\le s_0$. Summing yields
$\ell+\ell_1+\ell_2\le0$. All three lengths vanish, and
$s_1=s_2=s_0$. Thus $F$ is a bundle and its stable factors are
bundles of vector $v_0$. Extensions of locally free sheaves are
locally free, consistently with this conclusion.

For vector $v_0$, the same coprimality rules out an equal-slope
proper factor. If $A$ is such a stable sheaf, its reflexive hull
has square $q-2r_0\ell$, and \eqref{eq:stable-bound} again forces
$\ell=0$.

Slope stability implies Gieseker stability, and Gieseker
semistability implies slope semistability. A strictly
slope-semistable $F$ of vector $2v_0$ has the two factors just
described. Their reduced Hilbert polynomials both equal that of
$F$, so $F$ is Gieseker-semistable and is not Gieseker-stable.
For vector $v_0$ there is no strictly slope-semistable case.
The Jordan--H\"older factors in the two theories consequently
coincide, proving the S-equivalence assertion.
\end{proof}

\subsection{Geometry of the moduli space}
The vector $v_0$ is primitive. On a surface of Picard rank one,
$H$ is general in the sense required for these moduli spaces
\cite[Definition~2.8 and Lemma~2.9(1)]{pr23}. The primitive
space $N$ is a nonempty smooth projective irreducible
holomorphic symplectic manifold of dimension
\begin{equation}\label{eq:n}
 n=q+2,
\end{equation}
deformation equivalent to $S^{[q/2+1]}$
\cite[Theorem~1.7 and the paragraph following Definition~1.6]{pr23}.
In particular, it is simply connected. Proposition~\ref{prop:no-bubbling}
identifies it with a compact moduli space of slope-stable bundles.

The results of Kaledin--Lehn--Sorger
\cite[Theorem~4.4, Proposition~6.1 and Theorems~A--B]{kls06}
give the following description. The variety $M$ is normal,
irreducible and projective, of dimension $4q+2$. Its regular
locus is precisely the nonempty stable locus. Every singular
point is represented by $G_1\oplus G_2$ with $G_i\in N$.
The direct-sum morphism identifies the reduced singular locus
with $\Sym^2N$. One can also see the last assertion from the
local quiver descriptions below: the semisimple nonsimple
representations are exactly the unordered pairs of simple
summands. Near a repeated summand, the restrictions of traces
of matrix words to two one-dimensional characters are the
polarized power sums. These generate the symmetric-square
invariant ring in characteristic zero, identifying the
reduced singular locus. Away from the diagonal the local
product description gives the same identification. Accordingly,
\begin{equation}\label{eq:strata}
 \Sing M=\Sigma=\Sym^2N,\qquad
 \Delta\cong N,\qquad
 \dim\Sigma=2q+4,\quad \dim\Delta=q+2.
\end{equation}
There are three smooth strata: $M^s$, $\Sigma\setminus\Delta$,
and $\Delta$.

The variety $M$ has symplectic singularities and is locally
factorial in the algebraic sense. Its singular locus has
codimension $2q-2\ge6$, so its singularities are terminal
\cite[Corollary~1]{namikawa01}. It has no proper symplectic
resolution \cite[Theorem~B]{kls06}. The Main Theorem of
\cite{namikawa06} applies to the projective, $\mathbb Q$-factorial,
terminal symplectic variety $M$ and makes every flat deformation
locally trivial. Since $\Sigma$ is nonempty, $M$ has no flat
smoothing. Algebraic local factoriality here does not assert
factoriality of completed local rings; see
\cite[Remark~6.3]{kls06}.

We shall also use the irreducible symplectic-variety theorem of
Perego--Rapagnetta \cite[Theorem~1.10]{pr23}. It gives canonical
singularities and
\begin{equation}\label{eq:forms-cohomology}
 H^0(M,\Omega_M^{[2]})=\CC\sigma_M,
 \qquad H^1(M,\OO_M)=0.
\end{equation}
For the second equality, the definition of an irreducible
symplectic variety gives no reflexive one-forms. Pullback to
a projective resolution, Hodge symmetry, and rationality of
canonical singularities give the stated cohomology vanishing;
compare \cite[Definition~1.2 and Remark~1.3]{pr23}.
Symplectic singularities are Gorenstein, with trivial canonical
bundle generated by the top exterior power of $\sigma_M$.

\subsection{The two extension cones}\label{subsec:quivers}
The formulas below use the harmonic Hermitian forms; the displayed
complex vector spaces are supplied with the corresponding flat
hyperk\"ahler structures. A quiver representation is called
simple if it has no nonzero proper subrepresentation. Let
$G_1\not\cong G_2$ belong to $N$.
Stability and Serre duality give
\[
 \Hom(G_1,G_2)=\Ext^2(G_1,G_2)=0,
 \qquad \dim\Ext^1(G_1,G_2)=q.
\]
Each self-extension space has dimension $n$. The effective
unitary stabilizer is a circle, acting with opposite weights
on the two cross-extension spaces. Hence
\begin{equation}\label{eq:offdiag-cone}
 \Cone_{G_1\oplus G_2}
 =\CC^{2n}\times\mathcal N_q,
 \qquad \mathcal N_q=T^*\CC^q\hkr_0\U(1).
\end{equation}
In coordinates $(x,y)\in\CC^q\oplus(\CC^q)^*$ the equations are
$y(x)=0$ and $|x|^2=|y|^2$. The map $(x,y)\mapsto x\otimes y$
identifies the complex quotient with the rank-at-most-one,
square-zero matrices in $\mathfrak{sl}_q$, namely the minimal
nilpotent orbit closure. Its complex dimension is $2q-2$.
There are simple balanced points: take nonzero perpendicular
vectors of equal norm for $x$ and the conjugate of $y$.

For $E=G\oplus G$, the extension space is
$\Ext^1(G,G)\otimes\mathfrak{gl}_2$. Splitting off scalar
matrices gives
\begin{equation}\label{eq:diagonal-cone}
 \Cone_{G\oplus G}=\CC^n\times\mathcal Q_q,
 \qquad
 \mathcal Q_q=(\CC^n\otimes\mathfrak{sl}_2)\hkr_0\SU(2),
 \qquad \dim\mathcal Q_q=3q.
\end{equation}
Here $n$ is even, and a symplectic basis identifies the extension
space with $n/2$ pairs of matrix loops $(A_j,B_j)$.
Simple balanced points exist: set all $B_j=0$ and choose two
noncommuting Hermitian traceless matrices among the $A_j$, with
the remaining matrices zero. The complex moment map
$\sum_j[A_j,B_j]$ and the real moment map
$\sum_j([A_j,A_j^*]+[B_j,B_j^*])$ vanish. The two Hermitian
matrices have no common invariant line, so this representation
is simple. Scaling makes these simple points arbitrarily small.
The scalar center acts trivially, so the use of $\SU(2)$ rather
than its effective quotient does not change the metric quotient.

The holomorphic germs and metric normalizations in
\eqref{eq:offdiag-cone}--\eqref{eq:diagonal-cone} will be identified
with those of $M$ in the next section. In particular, the metric
on the fixed scalar factor of \eqref{eq:diagonal-cone} is the
one induced by $a\mapsto a\otimes I_2$; it is twice the harmonic
metric on the deformation space of $G$.

\begin{example}\label{ex:quartic}
For a Picard-rank-one quartic K3 surface, $v_0=(5,H,0)$ satisfies
$q=4$, $r_0=q+1$ and $\gcd(5,4)=1$. The stable factors have
second Chern class $c_2=7$, and a sheaf of vector $2v_0$ has
$c_2=18$, since $s=r+c_1^2/2-c_2$ on a K3 surface. The
dimensions are as follows.
\begin{center}
\begin{adjustbox}{max width=\textwidth}
\begin{tabular}{lcc}
\toprule
Locus & Complex dimension & Complex normal-cone dimension\\
\midrule
$M^s$ & $18$ & $0$\\
$\Sigma\setminus\Delta$ & $12$ & $6$\\
$\Delta$ & $6$ & $12$\\
\bottomrule
\end{tabular}
\end{adjustbox}
\end{center}
Thus the family contains a compact example with both types of
singular points. The existence assertions are the moduli
theorems above; the numerical equalities alone do not establish
existence.
\end{example}

\section{Harmonic coordinates and the quadratic estimate}\label{sec:harmonic}

Let $E=\bigoplus_i F_i\otimes W_i$ be a polystable bundle on
$S$, with pairwise nonisomorphic stable factors and a direct-sum
Hermitian--Einstein metric $h_E$. Put
\[
 V=H^1(S,\End E),\qquad
 G=\prod_i\GL(W_i),\qquad K=\prod_i\U(W_i).
\]
We use unnormalized trace inner products and volume form
$\omega^2/2$. Write $d=\dbar_E$, and let $H_j,G_j$ denote
harmonic projection and the Green operator in degree $(0,j)$.
The flat metric on $V$ is the real metric of this harmonic
Hermitian form. Its quaternionic structure is induced by the
parallel holomorphic symplectic form of $S$, normalized to make
that structure antiunitary. All quotient metrics use these
same normalizations.

\subsection{Exact quadraticity in the harmonic coordinates}
The canonical harmonic recursion is
\begin{equation}\label{eq:recursion}
 a(x)=x-d^*G_2(a(x)\wedge a(x)),\qquad
 \kappa_E(x)=H_2(a(x)\wedge a(x)).
\end{equation}
It converges on a neighbourhood of $0\in V$ and gives a
holomorphic semiuniversal deformation over
$Z=\kappa_E^{-1}(0)$; its underlying reduced space suffices
for all metric statements. In particular,
\begin{equation}\label{eq:coclosed}
 H_1a=x,\qquad d^*a=0,\qquad a=x+O(|x|^2).
\end{equation}
The maps are locally $G$-equivariant. These assertions, including
integrability precisely on $Z$, are the harmonic construction
of \cite[Lemma~2.1]{wuebben15}. Convergence follows from the
contraction for $d^*G_2(a\wedge a)$ in a sufficiently high
Sobolev norm. Parallel endomorphisms commute with the Hodge
operators, which gives the equivariance even for nonunitary
elements of $G$ whenever both sides are defined.

\begin{lemma}\label{lem:quadraticity}
On the full convergence neighbourhood of \eqref{eq:recursion},
\begin{equation}\label{eq:exact-quadraticity}
 \kappa_E(x)=H_2(x\wedge x).
\end{equation}
\end{lemma}

\begin{proof}
Choose a nonzero parallel holomorphic two-form $\sigma$ on $S$.
Every endomorphism-valued $(0,2)$-form is $\phi\bar\sigma$.
The formula for $d^*$, together with parallelness of $\sigma$,
identifies $d^*(\phi\bar\sigma)$ with a nonzero constant
multiple of contraction of $\partial_E\phi$ by $\bar\sigma$.
Thus $\phi\bar\sigma$ is harmonic if and only if
$\partial_E\phi=0$. Equivalently, $\phi^*$ is holomorphic.
The induced connection on $\End E$ has zero Hermitian--Einstein
constant, so its holomorphic sections are parallel. We obtain
\begin{equation}\label{eq:harmonic-two}
 \mathcal H^{0,2}(\End E)
   =\{\xi\bar\sigma:\nabla\xi=0\}.
\end{equation}
Consequently, for $\theta=\theta_0\bar\sigma$,
\begin{equation}\label{eq:harmonic-test}
 H_2\theta=0\quad\Longleftrightarrow\quad
 \int_S\tr(\theta_0\xi^*)\,dV=0
 \quad\hbox{for every parallel }\xi.
\end{equation}

We claim that for $d^*u=0$ and arbitrary
$\gamma=\phi\bar\sigma$,
\begin{equation}\label{eq:coexact-vanishing}
 H_2(u\wedge d^*\gamma)=0
      =H_2(d^*\gamma\wedge u).
\end{equation}
Use unitary normal frames at a point. Up to fixed nonzero
scalars, the two coefficients relative to $\bar\sigma$ are
\[
 -\sum_j u_{\bar j}\nabla_j\phi,
 \qquad \sum_j(\nabla_j\phi)u_{\bar j}.
\]
For a parallel $\xi$, covariant integration by parts gives
\begin{align}
 \int_S\sum_j\tr(u_{\bar j}\nabla_j\phi\,\xi^*)\,dV
 &=-\int_S\tr\left(\left(\sum_j\nabla_j u_{\bar j}\right)
                   \phi\xi^*\right)dV=0,\label{eq:ordered-one}\\
 \int_S\sum_j\tr((\nabla_j\phi)u_{\bar j}\xi^*)\,dV
 &=-\int_S\tr\left(\phi\left(\sum_j\nabla_j u_{\bar j}\right)
                   \xi^*\right)dV=0.\label{eq:ordered-two}
\end{align}
The contracted derivative is a constant multiple of $-d^*u$.
The connection preserves the metric and volume, and
$\nabla\xi^*=0$, so these are global divergence identities.
The order of the endomorphisms is retained in both formulas;
neither commuting factors nor flatness of the Chern connection
is used. Equations \eqref{eq:harmonic-test}--\eqref{eq:ordered-two}
prove \eqref{eq:coexact-vanishing}.

Now write $a=x+\beta$, where
$\beta=-d^*G_2(a\wedge a)$. Both $x$ and $\beta$ are
$d^*$-closed. Formula \eqref{eq:coexact-vanishing}, first with
$u=x$ and then with $u=\beta$, annihilates under $H_2$ all
three terms $x\wedge\beta$, $\beta\wedge x$ and
$\beta\wedge\beta$. Expanding $H_2(a\wedge a)$ proves
\eqref{eq:exact-quadraticity}.
\end{proof}

This calculation specifies the coordinates in the quadraticity
theorems for hyperholomorphic bundles. It uses the central
Ricci-flat metric on the surface. The coordinate specification
is essential for the metric estimate below: a general
holomorphic substitution $x\mapsto x+O(|x|^2)$ has derivative
$I+O(|x|)$ and need only preserve a relative linear error.

Choose unitary harmonic bases for the cross-extension spaces
and symplectically paired bases for the loops. Serre duality
identifies $H_2(x\wedge x)$ with the complex moment map of this
doubled quiver. More precisely, with $r_i=\rk F_i$ and a
normalized harmonic generator $\eta_S$ of $H^2(S,\OO_S)$,
\begin{equation}\label{eq:rank-dual}
 H_2(x\wedge x)
   =\bigoplus_i r_i^{-1}I_{F_i}\otimes
                      \mu_{\CC,i}(x)\,\eta_S,
\end{equation}
up to the single common nonzero symplectic normalization;
see the rank-dual formula in \cite[Section~3]{wuebben14}.
In particular the zero sets agree exactly. Since
$H^1(S,\OO_S)=0$, fixing the determinant removes no harmonic
direction. Lemma~\ref{lem:quadraticity} therefore gives
\begin{equation}\label{eq:Z-cone}
 Z=\mu_{\CC}^{-1}(0),\qquad
 (Z\cap\mu_{\RR}^{-1}(0))/K=\Cone_E.
\end{equation}

\subsection{The quantitative local comparison}
We record precisely the metric input from
\cite[Theorem~1.1(i),(iii)]{wuebben15}, at constant polarization.
For a fixed locally free polystable background on a compact
K\"ahler manifold, in the coordinates \eqref{eq:recursion} and
on the fixed-determinant subspace, that theorem gives a
homeomorphism
\[
 f_0:(Z\cap\mu_{\RR}^{-1}(0))/K\longrightarrow
       \text{a neighbourhood of }[E]
\]
onto the nearby Hermitian--Einstein moduli. It covers every
sufficiently nearby solution, preserves the local complex-orbit
description, and is smooth on smooth pieces of fixed orbit type.
For genuine tangent vectors to such pieces,
\begin{equation}\label{eq:imported-estimate}
 (1-C|x|^2)g_{\rm fin}
 \le f_0^*g^{L^2}\le(1+C|x|^2)g_{\rm fin}.
\end{equation}
Here, if $v$ is a balanced tangent lift and $C_x\xi=[\xi,x]$,
\[
 g_{{\rm fin},x}([v],[v])=\inf_\xi\|v-C_x\xi\|^2
\]
is the flat reduction metric; the scalar kernel can be omitted
from the gauge algebra. This is a statement on actual stratum
tangents, not on all Zariski tangents at a singular point.
The constant is independent of stabilizer type, of the positive
singular values of $C_x$ or $D\kappa_E(x)$, and of ratios of
positive orbit eigenvalues. It may depend on the fixed
background, ranks and multiplicities, harmonic dimensions,
positive elliptic gaps, Sobolev constants and the chosen
regularity order. These are exactly the dependencies needed
in \eqref{eq:main-comparison}.

All hypotheses hold here by Proposition~\ref{prop:no-bubbling}.
The Hermitian--Einstein correspondence identifies the analytic
moduli with the slope-polystable bundle moduli, and the latter
is the Gieseker moduli germ by that proposition. By
\eqref{eq:Z-cone}, $g_{\rm fin}=g_{\Cone_E}$. A zero-level
balanced representative satisfies
\begin{equation}\label{eq:cone-radius}
 r([x])=\|x\|.
\end{equation}
Indeed, the radial segment is horizontal and has length
$\|x\|$, while the norm is $1$-Lipschitz on every stratum.
The full cone distance has the same property
\cite[Section~5]{wuebben14}. Thus \eqref{eq:imported-estimate}
is precisely the asserted quadratic tensor estimate.

It remains to check that the identification is holomorphic in
these same coordinates. The family $\dbar_E+a(x)$ is
holomorphic, locally $G$-equivariant and semiuniversal. The
analytic GIT construction of Buchdahl--Schumacher
\cite[Theorems~5--9 and Remark~4]{buchdahl21} identifies its
reduced analytic quotient $Z/\!/G$ with the reduced moduli
germ. Denote its holomorphic classifying isomorphism by $\chi$.
The Kempf--Ness identification supplies the homeomorphism
between the balanced quotient in \eqref{eq:Z-cone} and
$Z/\!/G$.

The map $f_0$ in \eqref{eq:imported-estimate} corrects the
same operator $\dbar_E+a(x)$ within its complex gauge orbit.
Its finite balancing step changes $x$ within its $G$-orbit.
Consequently $f_0$ and $\chi$ assign the same bundle class to
every balanced representative. Local coverage and uniqueness
identify them. The holomorphic structure on the cone quotient
therefore makes $\Psi_E=f_0$ a biholomorphism of germs, without
any further change of coordinates. This proves part~(ii) of
Theorem~\ref{thm:main}.

\subsection{Finite-dimensional slices used below}\label{subsec:slice-input}
We also need two finite-dimensional conclusions of
\cite[Sections~4--6]{wuebben14}. At a fixed point of a flat
zero-level doubled-quiver quotient, its full quaternionic
normal representation $W$ is again a doubled Ext-quiver
representation. If $K'$ is the stabilizer and $U=W^{K'}$,
the normal quotient splits as the exact metric product
\begin{equation}\label{eq:normal-product}
 W\hkr_0K'=U\times\bigl(U^\perp\hkr_0K'\bigr).
\end{equation}
The stratum through the centre is $U\times\{0\}$. There is a
slice homeomorphism, smooth on strata, with relative metric
error $O(\|w\|)$ and derivative at the central stratum equal
to the isometric normal identification. Its constant may
depend on that centre, but not on positive eigenvalues of the
retained orbit operator.

The quotient is proper and geodesic, with its quotient
topology. If its complex zero fibre has one module-simple
point, then the simple locus is dense and every full normal
Ext-quiver has a module-simple point as well. These are the
metric-slice theorem, the properness theorem and the
normal-simple lemma of \cite{wuebben14}. The explicit simple
points in Section~\ref{subsec:quivers} verify this hypothesis
for our two singular cones. These statements apply to the
full $W$ in \eqref{eq:normal-product}; the complementary
representation $U^\perp$ alone need not be a doubled quiver.

\section{Completion and tangent cones}\label{sec:distances}

We distinguish two path distances until their equality has
been proved. On $M_{\reg}$, let $d_{\reg}$ be the Riemannian
distance of $g^{L^2}$. On all of $M$, call a path admissible
if it is continuous, has a finite partition into smooth paths
lying in single strata on the open subintervals, and has
finite one-sided stratum lengths at the endpoints. Let
$d_{\str}$ be the infimum of these lengths. The same
convention on each finite cone gives its full length distance
$d_{\Cone_E}$. Piecewise smooth paths and their finite-length
limits give the same infimum.

\subsection{Short paths in the regular locus}
The following elementary consequence of semialgebraicity
provides control near the endpoints of transverse perturbations.

\begin{lemma}\label{lem:short-paths}
Let $\Cone$ be one of the cones in
\eqref{eq:offdiag-cone}--\eqref{eq:diagonal-cone}, or a full
normal cone to one of its strata. There is a constant $A_\Cone$
such that any two regular points of radius at most $\rho$
can be joined in the regular locus, within that radial ball,
by a piecewise smooth path of length at most $A_\Cone\rho$.
\end{lemma}

\begin{proof}
Write the full normal representation as $W$ with compact
group $K'$, and let $B^*$ be the module-simple part of the
simultaneous zero set of its real and complex moment maps.
The simple locus of $\Cone$ is a connected complex open
subset of an irreducible cone. For the two original cones,
irreducibility and normality follow from their moduli germs;
homogeneity extends these properties to the whole cone.
The same argument applies to full normal slices, using the
slice theorem. Simplicity agrees with regularity here by
the stable-locus description of Section~\ref{sec:moduli}
and its slice restriction. Existence and density of simple
points in every normal cone were recalled in
Section~\ref{subsec:slice-input}.

The map $B^*\to\Cone_{\reg}$ is proper with connected
$K'$-orbit fibres. The group $K'$ is a product of unitary
groups, with an optional ineffective scalar quotient, and
is connected. Thus $B^*$ is connected. If the zero
representation itself is simple, the effective group is
trivial and the quotient is a flat vector space; the
assertion is then immediate. Otherwise radial scaling
identifies $B^*$ with its unit section times $(0,\infty)$.
It follows that
\[
 B^*\cap\{0<\|w\|<1\}
\]
is a bounded connected semialgebraic set.

We recall why such a set has finite Euclidean inner diameter,
including when it is not closed. A finite semialgebraic
triangulation uses finitely many open simplices. In each
simplex choose an anchor and join a variable point to that
anchor along the image of the straight segment in the
simplex. Join the finitely many anchors by fixed
semialgebraic arcs in the connected set. This gives finitely
many semialgebraic families of paths. Uniform monotonicity
partitions each coordinate of each path into a bounded
number of monotone intervals. Each coordinate varies over
a bounded range, so its total variation, and hence the
Euclidean length, has a common bound. This is the
triangulation and monotonicity argument of
\cite[Sections~2 and~3, especially Theorem~3.4]{vdd98}.
Semialgebraic arcs are piecewise smooth; endpoint derivatives
may be unbounded, but their lengths are finite.

Projection from the balanced locus to its smooth quotient
does not increase lengths. We obtain a uniform bound for
regular points of cone radius less than one. Dilation gives
the bound $A_\Cone\rho$ for an open radial ball. For points
on its boundary, first contract both endpoints by a factor
$1-\delta$, join the contracted points, and append the two
regular radial segments. Their total added length is at
most $2\delta\rho$. This proves the assertion for the
closed ball as well, after increasing the constant.
\end{proof}

\begin{proposition}\label{prop:completion}
The completion of $(M_{\reg},d_{\reg})$ is canonically
homeomorphic to $M$, and its distance equals $d_{\str}$.
In particular it is a compact geodesic space.
\end{proposition}

\begin{proof}
Shrink each fixed chart in Theorem~\ref{thm:main}(ii) so
that its squared metric bounds lie between $1/2$ and $2$.
Lemma~\ref{lem:short-paths} then shows that the $d_{\reg}$
diameter of regular points in a radial ball of radius
$\rho$ is at most $A\rho$. Every sequence of regular
points converging to the centre therefore determines the
same completion point.

There is also a lower bound for leaving a chart. If an
admissible path starts at cone radius $\rho$ and first
reaches the boundary of a concentric chart of radius $R$,
its length before exit is at least
\begin{equation}\label{eq:exit}
 (R-\rho)/\sqrt2.
\end{equation}
This follows by integrating the radial $1$-Lipschitz
inequality for the cone and the coarse tensor lower bound.
In particular, distinct points of $M$ cannot represent
the same completion point. Conversely, a $d_{\reg}$-Cauchy
sequence has an algebraically convergent subsequence because
$M$ is compact. The diameter estimate identifies its limit
in the completion; the Cauchy property gives the same limit
for the full sequence. Thus no additional completion points
occur. These estimates also prove that the completion
topology is the original topology of $M$. Denote this
completed distance by $D$.

We next compare lengths. Let $\gamma:[a,b]\to M$ be a
$C^1$ path contained in a single smooth stratum. Given
$\eta>0$, cover its compact image by finitely many centred
charts whose metric distortion is between $1-\eta$ and
$1+\eta$. In each such chart the cone splits as
$U\times\mathcal N$, and the stratum through its centre
is $U\times\{0\}$. Subdivide $[a,b]$ so that each
subarc is contained in one chart and write it there as
$(u(t),0)$.

Choose regular points $z_j\in\mathcal N$ tending to
the vertex. They exist by the simple-point property in
Section~\ref{subsec:slice-input}. The paths
\begin{equation}\label{eq:transverse-path}
 t\longmapsto (u(t),z_j)
\end{equation}
lie in the regular locus for large $j$. Their lengths in
the exact product are the Euclidean length of $u$.
The two tensor bounds give, for each subarc,
\[
 \limsup_j L_{g^{L^2}}(\Psi(u,z_j))
 \le \sqrt{\frac{1+\eta}{1-\eta}}\,L_{g^{L^2}}(\gamma).
\]
Their endpoints converge in the original topology, hence
in $D$ by the preceding paragraph. Taking endpoint
distance limits and summing with the triangle inequality
gives the same upper bound for $D(\gamma(a),\gamma(b))$.
No common endpoint is required for the different
perturbed subarcs. Letting $\eta$ tend to zero gives
$D(\gamma(a),\gamma(b))\le L(\gamma)$.

For an admissible path whose open pieces have endpoints
in lower strata, truncate each piece at both ends and
apply this argument. The finite one-sided lengths and
continuity of $D$ allow the truncation parameters to
tend to zero. Therefore $D\le d_{\str}$ on $M$.
On regular points, every regular path is admissible, so
$d_{\str}\le d_{\reg}=D$. Finally, $d_{\str}$ is
continuous in the original topology: in a chart centred
at any point, its radial paths have length at most
$\sqrt2\rho$, and the triangle inequality gives
continuity in both arguments. Approximation by regular
endpoints extends the reverse inequality to all of $M$.
Thus $D=d_{\str}$.

The completion of a length space is a length space.
It is compact here, so the usual compactness argument
for almost minimizing paths supplies minimizing
geodesics. This proves the last assertion.
\end{proof}

\subsection{Quantitative convergence to the cone}
We can now transfer the quadratic tensor comparison to
the completed distance. Fix a centre $p=[E]$ and its
chart $\Psi_E$. Competitors joining points of radius
at most $\rho$ can be chosen with length $O(\rho)$,
using radial paths through the centre. A competitor
of this length cannot leave a fixed chart when $\rho$
is small, by \eqref{eq:exit}. Within the chart, the
coarse lower bound further shows that its radius is
$O(\rho)$ throughout. On all such competitors the
tensor distortion is therefore $1+O_E(\rho^2)$.

Apply this reasoning first to almost minimizing cone
paths and then to almost minimizing admissible paths
in $M$. Proposition~\ref{prop:completion} identifies
the latter infimum with $D$. Taking both infima gives
\begin{equation}\label{eq:distance-comparison}
 (1-C'_E\rho^2)d_{\Cone_E}(x,y)
 \le D(\Psi_E x,\Psi_E y)
 \le(1+C'_E\rho^2)d_{\Cone_E}(x,y)
 \quad(r(x),r(y)\le\rho).
\end{equation}
The slightly larger radius used by the competitors
is absorbed into $C'_E$. Setting $y=0$ proves
\eqref{eq:radius-main}.

For completeness, fix $L<\infty$ and put
$x=\epsilon\hat x$, $y=\epsilon\hat y$ in
\eqref{eq:distance-comparison}. On the cone ball of
radius $L$, the map
$\hat x\mapsto\Psi_E(\epsilon\hat x)$ has distance
distortion $O_{E,L}(\epsilon^2)$ when the target
distance is divided by $\epsilon$. The radial formula
shows that the preimage of the closed target unit
ball lies between cone radii $1-O_E(\epsilon^2)$
and $1+O_E(\epsilon^2)$. Radial adjustment to radius
one costs $O_E(\epsilon^2)$ in the rescaled metric.
This gives a correspondence of the two closed unit
balls with that distortion, hence the
Gromov--Hausdorff estimate in
Corollary~\ref{cor:completion-main}. The construction
works for every sequence $\epsilon\downarrow0$;
the metric tangent cone is consequently unique.

\section{The determinant class and the canonical metric}\label{sec:canonical}

Write $R=2r_0$, $c=2aH$ and $h=c_1(H)=[\omega]$.
Let $\Omega$ be the K\"ahler form of the harmonic metric
on $M_{\reg}$. In the normalization of
Biswas--Schumacher, denote the Weil--Petersson form by
$\omega^{\WP}$. There is a fixed positive constant $b$
such that
\begin{equation}\label{eq:WP-normalization}
 \Omega=b\,\omega^{\WP}.
\end{equation}
This constant only converts trace and real-metric conventions;
the cone and actual metrics throughout the paper use the
same harmonic convention. We retain $b$ explicitly in
the class calculation.

\subsection{The positive determinant polarization}
Choose a sufficiently divisible positive integer $k$, a
divisor $D\in|H^k|$, and a point $x\in S$. For the fixed
Mukai vector $v=2v_0$ define
\begin{equation}\label{eq:determinant-class}
 u_D=-R[\OO_D]+\chi(F\otimes\OO_D)[\OO_x],
 \qquad
 \lambda_D=\det Rp_*(\mathcal E\otimes q_S^*u_D).
\end{equation}
The Euler characteristic depends only on $v$, not on $F$.
The expression on the right is the determinant of cohomology
on family parameter spaces, descended to $M$. It is the
positive convention of
\cite[Definition~8.1.9]{hl97}. Here $p$ and $q_S$ are the
projections to the parameter space and $S$. Any positive
power required in the construction can be absorbed into
$\lambda_D$, with the reciprocal change in the scalar
class coefficient below. The auxiliary $k$ does not change
the fixed form $\omega$ used for the harmonic metric.

The stabilizer weights also make descent explicit in this
family. For a stable factor $G$ of vector $v_0$, additivity
of the Euler characteristic gives
\[
 \chi(G\otimes u_D)
 =-2r_0\chi(G\otimes\OO_D)
             +2\chi(G\otimes\OO_D)r_0=0.
\]
Each multiplicity-group determinant character is therefore
trivial at a closed polystable orbit, including a repeated
summand. At a stable point the same equality follows directly
from the definition of $u_D$.

\begin{lemma}\label{lem:ample}
The determinant line $\lambda_D$ is ample on $M$.
\end{lemma}

\begin{proof}
The determinant construction is semiample and defines
the projective Donaldson--Uhlenbeck morphism
\cite[Definition and Theorem~8.2.8]{hl97}. Its closed-point
invariant is the reflexive hull of the slope graded object
together with the defect cycle of that graded object
\cite[Theorem~8.2.11]{hl97}. Proposition~\ref{prop:no-bubbling}
says that every slope graded factor here is locally free.
Thus the defect cycle is zero. The remaining invariant
is the polystable bundle itself, and slope S-equivalence
agrees with Gieseker S-equivalence. The morphism is
therefore injective on closed points. Being projective
and quasi-finite, it is finite. A positive power of
$\lambda_D$ is the pullback of an ample line under this
finite morphism, proving ampleness. This argument needs
neither normality of the target nor an isomorphism with it.
\end{proof}

\begin{lemma}\label{lem:class}
With the determinant convention \eqref{eq:determinant-class},
before replacing it by a descent power,
\begin{equation}\label{eq:omega-class}
 [\Omega]=\frac{4\pi^2b}{Rk}\,c_1(\lambda_D)
       \quad\hbox{in }H^2(M_{\reg},\RR).
\end{equation}
\end{lemma}

\begin{proof}
The Chern character of the $K$-theory class in
\eqref{eq:determinant-class} is
\begin{equation}\label{eq:chu}
 \ch(u_D)=(0,-Rkh,k(c\cdot h)[x]).
\end{equation}
Indeed, $\ch(\OO_D)=1-e^{-kh}$ and
$\chi(F\otimes\OO_D)=k(c\cdot h)-Rk^2h^2/2$
on a K3 surface.

Work first on a stable parameter space with a universal
family $\mathcal E$. Since $H^1(S,\RR)=0$,
\[
 c_1(\mathcal E)=q_S^*c+p^*\beta
\]
for a class $\beta$ on that space. Grothendieck--Riemann--Roch
and \eqref{eq:chu} give
\begin{align}
 c_1(\lambda_D)
 &=-Rk\,p_*(\ch_2(\mathcal E)\,q_S^*h)
                  +k(c\cdot h)\beta \notag\\
 &=-\frac{k}{2}\,p_*(\ch_2(\End\mathcal E)\,q_S^*h).
                                                    \label{eq:GRR}
\end{align}
For the second equality use
$\ch_2(\End\mathcal E)=2R\ch_2(\mathcal E)-c_1(\mathcal E)^2$
and
$p_*(c_1(\mathcal E)^2q_S^*h)=2(c\cdot h)\beta$.
This agrees with the discriminant formula in
\cite[Proposition~8.3.1(ii)]{hl97}.

The fixed-determinant curvature identity of
\cite[Proposition~4.3, equation~(4.6)]{biswas16}, in
surface dimension, is
\begin{equation}\label{eq:WP-curvature}
 \frac{\omega^{\WP}}{2\pi^2}
   =-\frac1R\,p_*(\ch_2(\End\mathcal E,h_E)\wedge\omega).
\end{equation}
The fixed-determinant hypothesis holds because
$\operatorname{Pic}^0(S)=0$. Comparing
\eqref{eq:GRR} and \eqref{eq:WP-curvature} gives the stated
coefficient.

One can check the sign through the Quillen line as well.
In the convention of \cite[Theorem~4.6, equations~(4.11)
and~(4.13)]{biswas16}, the positive line is the dual
determinant
\begin{equation}\label{eq:Quillen}
 Q_D=\left[\det Rp_*\left(
 (\End\mathcal E-R^2\OO)\otimes
                   (q_S^*H^k-q_S^*H^{-k})\right)\right]^{-1}.
\end{equation}
Its first Chern class is
\[
 c_1(Q_D)=-2k\,p_*(\ch_2(\End\mathcal E)\,q_S^*h)
                  =4c_1(\lambda_D).
\]
Using the nondual determinant would reverse this sign.

Finally these computations give equality on the coarse
stable locus, not merely on a family. The stable Quot
parameter space is a principal $\operatorname{PGL}(m)$
bundle over $M^s$ for an appropriate framing dimension
$m$. Its universal endomorphism bundle and determinant
classes descend. Pullback in real degree-two cohomology
is injective, by the low-degree Leray sequence and
$H^1(\operatorname{PGL}(m),\RR)=0$. Thus the equality
proved on that atlas gives \eqref{eq:omega-class} on
$M^s=M_{\reg}$. No universal bundle on the full coarse
moduli space is required.
\end{proof}

Put $\lambda=\lambda_D$ and
$\kappa=4\pi^2b/(Rk)$, adjusted if a positive descent
power has been taken. Choose a smooth K\"ahler
representative $\theta\in c_1(\lambda)$, obtained from
a projective embedding and scaled by its degree.

\subsection{Continuous potentials in the global class}
Theorem~11 of \cite{buchdahl21} applies to polystable
bundles deformable to stable bundles over a compact Hodge
manifold. The surface $S$ is Hodge since $[\omega]=c_1(H)$;
all our polystable objects are bundles; and the simple
points in Section~\ref{subsec:quivers} show local density
of stable bundles. The analytic moduli germs are those
identified in Section~\ref{sec:harmonic}. That theorem
therefore gives continuous plurisubharmonic local
potentials for $\Omega$. We now check that they define
a current in the global class from Lemma~\ref{lem:class}.

The variety $M$ is Cohen--Macaulay and
$\operatorname{codim}\Sing M\ge6$. Scheja's restriction
theorem \cite[Section~5, Corollary after Satz~III]{scheja61}
gives
\begin{equation}\label{eq:Scheja}
 H^1(M,\OO_M)\simeq H^1(M_{\reg},\OO_{M_{\reg}}).
\end{equation}
In the theorem's depth convention, for degree one it
suffices that
$\operatorname{depth}_x\OO_M\ge\dim_x\Sing M+3$;
our dimension and codimension bounds satisfy this.
Combining \eqref{eq:Scheja} with
\eqref{eq:forms-cohomology} gives vanishing of the
right side.

By Lemma~\ref{lem:class}, there is a smooth real
one-form $\alpha$ on $M_{\reg}$ with
$d\alpha=\Omega-\kappa\theta$. Its $(0,1)$ part is
$\dbar$-closed. The preceding vanishing and the
Dolbeault theorem give a smooth function $f$ with
$\alpha^{0,1}=\dbar f$. Since $\alpha$ is real,
\begin{equation}\label{eq:global-potential}
 \Omega=\kappa\theta+dd^c\varphi
 \quad\hbox{on }M_{\reg},\qquad
 \varphi=2\operatorname{Im}f.
\end{equation}

Let $U$ be a small neighbourhood in $M$. Write
$\Omega=dd^cu$ on $U_{\reg}$, where $u$ is the
continuous local potential supplied above, and
$\theta=dd^cw$ with $w$ a smooth local potential.
Then
\[
 h_U=\varphi-u+\kappa w
\]
is pluriharmonic on $U_{\reg}$. Both $h_U$ and
$-h_U$ are plurisubharmonic there. The
Grauert--Remmert extension theorem across a subset of
complex codimension at least two in a normal space
\cite[Section~1.7, Satz~4]{gr56} extends both as
plurisubharmonic functions on $U$. In this codimension
the theorem requires no prior upper bound. Their sum
is zero on the dense regular locus, hence everywhere
by uniqueness of plurisubharmonic extension. The
extensions are consequently finite and continuous:
each is both upper and lower semicontinuous. Their
$dd^c$ currents are opposite positive currents, so
both vanish.

It follows that $\varphi=u-\kappa w+h_U$ extends
continuously over $U$. The extensions agree on overlaps
by density, producing a global continuous function
$\varphi$ on $M$. Formula \eqref{eq:global-potential}
now defines a positive closed current on $M$ in
$\kappa c_1(\lambda)$, with the claimed local
continuous potentials.

\subsection{Identification by Ricci-flat uniqueness}
The harmonic metric on the stable locus is
hyperk\"ahler. This is Itoh's theorem for irreducible
Einstein--Hermitian connections on a K3 surface
\cite[Theorem~2.1 and Remark~(iii)]{itoh88}; the
Einstein--Hermitian extension in that remark allows
the nonzero slope here. The associated holomorphic
symplectic form, in harmonic coordinates, is a fixed
nonzero multiple of
\begin{equation}\label{eq:Mukai-form}
 (a,b)\longmapsto\int_S\tr(a\wedge b)\wedge\sigma.
\end{equation}
Indeed, the quaternionic complex structures on
endomorphism-valued one-forms and the trace $L^2$
pairing give \eqref{eq:Mukai-form} by contraction
with the parallel form $\sigma$. This is the
Serre-duality form defining the Mukai symplectic
structure, so it agrees with $\sigma_M$ up to the
fixed normalization. This identification is a
calculation of the forms; it does not require
extending a Riemannian tensor across $\Sigma$.

Set $d=\dim_{\CC}M=4q+2$. The hyperk\"ahler volume
identity on $M_{\reg}$ says that $\Omega^d$ is a
constant positive multiple of the adapted measure
defined by $\sigma_M^{d/2}$ and its conjugate. The
constant is determined by the total volume of
$\kappa c_1(\lambda)$. The current has locally
bounded potentials, so its Monge--Amp\`ere measure
does not charge the singular analytic set. The
volume identity therefore holds on all of $M$.

The variety is projective, has canonical
singularities and trivial canonical bundle. The
class $\kappa c_1(\lambda)$ is K\"ahler and belongs
to the real N\'eron--Severi space. The existence and
uniqueness theorem of Eyssidieux--Guedj--Zeriahi
\cite[Theorem~7.5 and Corollary~7.6]{egz09} applies
to precisely these data. It identifies our current
with the canonical singular Ricci-flat current in
that class. This proves part~(i) of
Theorem~\ref{thm:main}. Together with
Section~\ref{sec:distances}, it proves all claims
of Corollary~\ref{cor:completion-main} for the
canonical metric.

\section{The metric on the singular strata}\label{sec:strata}

We finish by distinguishing the quadratic approximation
from an exact local cone metric. The argument uses only
the intrinsic recognition of the two singular strata and
their harmonic metrics.

\begin{lemma}\label{lem:intrinsic-strata}
The subset $\Sigma$ consists precisely of the points of
$M$ whose metric tangent cones are not Euclidean. Within
$\Sigma$, the subset $\Delta$ is its topological
nonmanifold locus. The corresponding descriptions hold
in the cones \eqref{eq:offdiag-cone} and
\eqref{eq:diagonal-cone}.
\end{lemma}

\begin{proof}
At a regular point the tangent cone is Euclidean.
At an off-diagonal point the normal cone is
$\mathcal N_q=T^*\CC^q\hkr_0\U(1)$. Its link is
\begin{equation}\label{eq:link}
 L_q=V_2(\CC^q)/S^1,
\end{equation}
where $V_2(\CC^q)$ is the complex Stiefel manifold
of orthonormal two-frames and the circle multiplies
both vectors by the same phase. To see this, conjugate
the cotangent coordinate in the equations
$y(x)=0$, $|x|=|y|$ and then fix the common norm.

The fibration
$S^{2q-3}\to V_2(\CC^q)\to S^{2q-1}$ gives
$\pi_1(V_2)=\pi_2(V_2)=0$ for $q\ge4$.
The circle bundle over $L_q$ consequently gives
$\pi_1(L_q)=0$ and $\pi_2(L_q)=\ZZ$. Hurewicz
then gives $H_2(L_q,\ZZ)=\ZZ$. Thus the vertex
of $\mathcal N_q$ has nonzero local homology in
degree three, strictly below its real dimension
$4q-4$. Multiplication by the Euclidean factor
shifts this nonzero local homology but does not
remove it. The off-diagonal cone is therefore not
a topological manifold at its vertex, and in
particular is not Euclidean.

The diagonal cone contains off-diagonal points
arbitrarily close to its vertex: choose two
distinct one-dimensional loop characters tending
to zero and take their direct sum. At each such
point the normal slice has the off-diagonal form
just considered. Hence the diagonal cone also
cannot be Euclidean. Uniqueness of tangent cones
from Section~\ref{sec:distances} proves the first
assertion on $M$. The finite-dimensional slice
theorem gives the same assertion at all points
of each cone.

By \eqref{eq:strata}, $\Sigma=\Sym^2N$.
Off the diagonal it is a manifold. Near a
diagonal point its germ is
\begin{equation}\label{eq:symmetric-square}
 \CC^n\times(\CC^n/\{\pm1\}).
\end{equation}
The normal link is $\mathbb{RP}^{2n-1}$, whose
first integral homology is $\ZZ/2$. This gives
nontop local homology at the vertex, so the
nonmanifold locus of $\Sigma$ is precisely
$\Delta$. In the diagonal cone, the nonsimple
semisimple representations are unordered pairs
of one-dimensional loop characters. Their
space is $\Sym^2\CC^n$, with the same description
\eqref{eq:symmetric-square}; its nonmanifold
locus is the diagonal flat factor $\CC^n$.
In the off-diagonal cone the singular set is
the flat factor $\CC^{2n}\times\{0\}$ and is
everywhere a manifold. This proves all the
claimed intrinsic descriptions.
\end{proof}

\begin{lemma}\label{lem:stratum-length}
Let $S_0$ be either $\Delta$ or
$\Sigma\setminus\Delta$. The intrinsic length
metric obtained by restricting curves to $S_0$
and measuring their lengths with the completed
ambient distance is the Riemannian length metric
of its harmonic stratum metric. The same assertion
holds for the corresponding flat strata in the
tangent cones.
\end{lemma}

\begin{proof}
Fix $s\in S_0$. In the chart centred at $s$, the
stratum is the fixed factor $U\times\{0\}$.
The coarse tensor bounds and the exit estimate
make the restricted ambient distance locally
bi-Lipschitz equivalent to a smooth metric on
$S_0$. Thus a rectifiable curve has a locally
absolutely continuous parametrization in its
smooth coordinates.

For a differentiable curve $\gamma$ with
$\gamma(0)=s$, write its fixed-factor coordinate
as $u(t)=t\dot u(0)+o(t)$. The radius formula
\eqref{eq:radius-main} in this fixed centred
chart gives
\begin{equation}\label{eq:metric-speed}
 \lim_{t\to0}\frac{D(\gamma(t),\gamma(0))}{|t|}
       =\|\dot u(0)\|_U
       =\|\dot\gamma(0)\|_{g_{S_0}}.
\end{equation}
The last equality follows from the harmonic
normalization at the centre. The argument also
applies at almost every differentiability point
of an absolutely continuous curve. It is
pointwise: the comparison constant belongs to
the fixed centre at which the limit is taken.

The length of a Lipschitz curve in a metric
space is the integral of its metric speed
\cite[Theorem~2.7.6]{bbi01}. Applying this
locally, and subdividing a rectifiable curve,
proves equality of its ambient-metric length
and harmonic Riemannian length. Infima over
curves in the stratum prove the assertion.
For the finite cones the same proof uses the
finite metric slice and its radius comparison.
\end{proof}

\begin{proposition}\label{prop:not-isometric}
For every singular point $p\in M$, no pointed
neighbourhood of $(M,D,p)$ is isometric to a
neighbourhood of the vertex in its metric
tangent cone.
\end{proposition}

\begin{proof}
An isometry preserves metric tangent cones,
and hence preserves $\Sigma$ by
Lemma~\ref{lem:intrinsic-strata}. If $p$ is
off the diagonal, a sufficiently small
neighbourhood meets only the smooth stratum
$\Sigma\setminus\Delta$. In the tangent
cone this stratum is the flat factor
$\CC^{2n}$. If $p$ is diagonal, an isometry
also preserves the topological nonmanifold
locus of $\Sigma$, so it sends $\Delta$
to the flat factor $\CC^n$ of the diagonal
cone. Lemma~\ref{lem:stratum-length} shows
that in either case the restricted map
would be an isometry of the intrinsic
Riemannian metrics on open subsets of
these strata.

On the local ordered cover of
$\Sigma\setminus\Delta$, the harmonic metric
is $g_N\oplus g_N$. On $\Delta$ it is
$2g_N$, since the tangent inclusion is
$a\mapsto a\otimes I_2$ and
$\tr(I_2)=2$. An isometry to a flat factor
therefore makes $g_N$ flat on a nonempty
open subset of $N$.

The metric $g_N$ is hyperk\"ahler, hence
Ricci-flat and real analytic in harmonic
coordinates \cite[Theorem~5.2]{dk81}.
Analytic continuation of its curvature on
the connected manifold $N$ would make it
flat everywhere. A complete simply
connected flat manifold is Euclidean,
whereas $N$ is compact, simply connected
and of positive dimension. This
contradiction proves the proposition.
\end{proof}

The proposition concerns exact isometry.
It provides no lower bound on the first
nonzero error term in
\eqref{eq:main-comparison}. Optimality of
the exponent, derivative estimates near
lower strata, and extensions to moduli
with a singular-sheaf boundary remain
separate questions.

\begin{thebibliography}{Wue26b}

\bibitem[AS18]{as18}
E.~Arbarello and G.~Sacc\`a,
\emph{Singularities of moduli spaces of sheaves on K3 surfaces and
Nakajima quiver varieties}, Adv. Math. \textbf{329} (2018), 649--703.
\href{https://arxiv.org/abs/1505.00759}{arXiv:1505.00759}.

\bibitem[BS16]{biswas16}
I.~Biswas and G.~Schumacher,
\emph{The Weil--Petersson current for moduli of vector bundles and
applications to orbifolds}, Ann. Fac. Sci. Toulouse Math. (6)
\textbf{25} (2016), no.~4, 895--917.
\href{https://doi.org/10.5802/afst.1514}{doi:10.5802/afst.1514}.

\bibitem[BSc21]{buchdahl21}
N.~Buchdahl and G.~Schumacher,
\emph{An analytic application of geometric invariant theory},
J. Geom. Phys. \textbf{165} (2021), 104237.
\href{https://doi.org/10.1016/j.geomphys.2021.104237}{doi:10.1016/j.geomphys.2021.104237}.
\href{https://arxiv.org/abs/2008.04625v3}{arXiv:2008.04625v3}.

\bibitem[BZ19]{bz19}
N.~Budur and Z.~Zhang, \emph{Formality conjecture for K3 surfaces},
Compos. Math. \textbf{155} (2019), no.~5, 902--911.
\href{https://arxiv.org/abs/1803.03974}{arXiv:1803.03974}.

\bibitem[BBI01]{bbi01}
D.~Burago, Y.~Burago and S.~Ivanov,
\emph{A Course in Metric Geometry}, Graduate Studies in Mathematics
\textbf{33}, American Mathematical Society, Providence, RI, 2001.

\bibitem[DK81]{dk81}
D.~M.~DeTurck and J.~L.~Kazdan,
\emph{Some regularity theorems in Riemannian geometry},
Ann. Sci. \'Ecole Norm. Sup. (4) \textbf{14} (1981), no.~3, 249--260.
\href{https://doi.org/10.24033/asens.1405}{doi:10.24033/asens.1405}.

\bibitem[EGZ09]{egz09}
P.~Eyssidieux, V.~Guedj and A.~Zeriahi,
\emph{Singular K\"ahler--Einstein metrics}, J. Amer. Math. Soc.
\textbf{22} (2009), no.~3, 607--639.
\href{https://doi.org/10.1090/S0894-0347-09-00629-8}{doi:10.1090/S0894-0347-09-00629-8}.
\href{https://arxiv.org/abs/math/0603431v2}{arXiv:math/0603431v2}.

\bibitem[GR56]{gr56}
H.~Grauert and R.~Remmert,
\emph{Plurisubharmonische Funktionen in komplexen R\"aumen},
Math. Z. \textbf{65} (1956), 175--194.
\href{https://doi.org/10.1007/BF01473877}{doi:10.1007/BF01473877}.

\bibitem[HL97]{hl97}
D.~Huybrechts and M.~Lehn,
\emph{The Geometry of Moduli Spaces of Sheaves},
Aspects of Mathematics E31, Vieweg, Braunschweig, 1997.
\href{https://doi.org/10.1007/978-3-663-11624-0}{doi:10.1007/978-3-663-11624-0}.

\bibitem[Ito88]{itoh88}
M.~Itoh, \emph{Gauge fields and quaternion structure},
Kodai Math. J. \textbf{11} (1988), no.~3, 344--360.
\href{https://doi.org/10.2996/kmj/1138038932}{doi:10.2996/kmj/1138038932}.

\bibitem[KLS06]{kls06}
D.~Kaledin, M.~Lehn and C.~Sorger,
\emph{Singular symplectic moduli spaces},
Invent. Math. \textbf{164} (2006), no.~3, 591--614.
\href{https://doi.org/10.1007/s00222-005-0484-6}{doi:10.1007/s00222-005-0484-6}.

\bibitem[Nam01]{namikawa01}
Y.~Namikawa, \emph{A note on symplectic singularities},
preprint, 2001.
\href{https://arxiv.org/abs/math/0101028}{arXiv:math/0101028}.

\bibitem[Nam06]{namikawa06}
Y.~Namikawa, \emph{On deformations of $\mathbb Q$-factorial
symplectic varieties}, J. Reine Angew. Math. \textbf{599} (2006), 97--110.
\href{https://doi.org/10.1515/CRELLE.2006.079}{doi:10.1515/CRELLE.2006.079}.

\bibitem[NO26]{no26}
Y.~Namikawa and Y.~Odaka,
\emph{Canonical torus action on symplectic singularities},
preprint, 2026.
\href{https://arxiv.org/abs/2503.15791v4}{arXiv:2503.15791v4}.

\bibitem[Oda26]{odaka26}
Y.~Odaka, \emph{Symplectic singularity as hyperK\"ahler cone},
preprint, 2026.
\href{https://arxiv.org/abs/2610.35510v1}{arXiv:2610.35510v1}.

\bibitem[PR23]{pr23}
A.~Perego and A.~Rapagnetta,
\emph{Irreducible symplectic varieties from moduli spaces of sheaves
on K3 and Abelian surfaces}, Algebr. Geom. \textbf{10} (2023),
no.~3, 348--393.
\href{https://doi.org/10.14231/AG-2023-012}{doi:10.14231/AG-2023-012}.
\href{https://arxiv.org/abs/1802.01182v2}{arXiv:1802.01182v2}.

\bibitem[Sch61]{scheja61}
G.~Scheja, \emph{Riemannsche Hebbarkeitss\"atze f\"ur
Cohomologieklassen}, Math. Ann. \textbf{144} (1961), 345--360.
\href{https://doi.org/10.1007/BF01470506}{doi:10.1007/BF01470506}.

\bibitem[vD98]{vdd98}
L.~van den Dries, \emph{O-minimal structures and real analytic geometry},
in \emph{Current Developments in Mathematics 1998},
International Press, Somerville, MA, 1999, pp.~105--152.

\bibitem[Ver96]{verbitsky96}
M.~Verbitsky, \emph{Hyperholomorphic bundles over a hyper-K\"ahler
manifold}, J. Algebraic Geom. \textbf{5} (1996), no.~4, 633--669.
\href{https://arxiv.org/abs/alg-geom/9307008v1}{arXiv:alg-geom/9307008v1}.

\bibitem[Wue26a]{wuebben14}
B.~J.~Wuebben, \emph{Quantitative metric wall crossing for K3 bundle
moduli}, preprint, 6 November 2026.
\href{https://doi.org/10.5281/zenodo.24024998}{doi:10.5281/zenodo.24024998}.

\bibitem[Wue26b]{wuebben15}
B.~J.~Wuebben, \emph{Quantitative local models for Hermitian--Einstein
metrics}, preprint, 7 November 2026.
\href{https://doi.org/10.5281/zenodo.24029090}{doi:10.5281/zenodo.24029090}.

\end{thebibliography}

\end{document}
